% ============================================================================
\section{Related Work}
\label{sec:related}

\subsection{Retrieval-augmented generation and reranking}
Retrieval-augmented generation conditions a generator on retrieved evidence
\cite{b1}, and dense bi-encoders such as Sentence-BERT made semantic retrieval
practical \cite{b2,b28}; Facebook AI Similarity Search (FAISS) provides the
nearest-neighbour infrastructure at scale \cite{b3}. Sparse retrieval remains a
strong baseline \cite{b4}, hybrid fusion by reciprocal rank fusion (RRF) is a
common robustness measure \cite{b5}, and cross-encoder reranking improves the top
of the ranking at higher cost \cite{b6}. Dense passage retrieval established the
two-stage recipe we follow \cite{b7}. All of this machinery is fixed in our study;
the object of study is the feedback signal applied on top of it. Reranking matters
most where the first stage is imperfect: a cross-encoder sees the query and
candidate jointly and can correct ordering errors that a bi-encoder cannot
\cite{b6}. RAG evaluation, in turn, separates retrieval quality from answer
quality \cite{b1,b40,b41}, a distinction that our proxy-versus-generated
divergence reproduces.

\subsection{Learning to rank from feedback signals}
Relevance feedback is one of the oldest ideas in information retrieval:
Rocchio-style query modification moves the query toward relevant documents
\cite{b8}, and relevance-based language models estimate a query model from feedback
documents \cite{b9}. Learning-to-rank inherits the same assumption---that logged
interactions carry usable relevance signal. Two questions that this literature
often leaves implicit are the ones we make explicit: \emph{what the feedback source
knows}, and \emph{how its evidence is converted into a score}. Feedback in
recommender systems, for instance, is known to carry popularity and base-rate bias
\cite{b10}; we show an analogous pathology in judge-driven ticket feedback,
identify its arithmetic cause, and repair it without labels. A second strand studies
implicit feedback and click models, where exposure and position bias must be
corrected before interactions can be treated as relevance \cite{b34,b37}. Our
setting removes exposure bias by judging every candidate in a pool, but it
introduces a different pathology: the judge's evidence is sparse and zero-heavy,
and the score construction amplifies that sparsity.

\subsection{LLM-as-judge validity}
Large language models are increasingly used as relevance and quality judges
\cite{b11,b12,b35,b36}. Their known failure modes include position bias, verbosity
bias, and self-preference \cite{b13,b14}. Our protocols differ from ordinary LLM
judging in one decisive respect: the conditioned judge receives the historical
resolution of the training query, which makes it a \emph{resolution-informed}
judge and an upper bound on what resolution-aware evidence can buy. It is
deployable for historical pairs: both tickets in a historical pair have known
resolutions, so the conditioned votes can be precomputed for the historical base.
At inference on a new ticket, however, no resolution exists yet, so only blind
votes can be generated for that ticket. We therefore report both protocols side by side, state exactly
what each judge and metric sees, and add an expert assessment of the generated
answers (Section~\ref{sec:expert}). Judge agreement is usually reported in
aggregate; the earlier evaluation found low per-ticket agreement between its
independent judge and the automatic metric, which we treat as a limitation of the
measurement rather than as evidence against the aggregate signal.

\subsection{Calibration, selective prediction, and abstention}
Converting scores into calibrated quantities is a classical problem
\cite{b15,b16,b17}, and empirical Bayes provides a standard way to stabilize rate
estimates built from few observations \cite{b18,b19,b38,b39}. Decision theory
prescribes abstention when expected utility is negative \cite{b20}; selective
prediction operationalizes this for classifiers \cite{b21}, and cost-sensitive
learning makes the trade-off explicit \cite{b22}. Our gate is a selective predictor
over tickets. Bayesian smoothing with a Beta or Dirichlet prior is standard in
language modelling and information retrieval \cite{b18,b19}; our empirical-Bayes
lift is a direct application of that idea to a per-candidate rate estimated from
very few votes. We do not claim the smoother as a methodological contribution; the
contribution is the demonstration that it removes harm without creating a confirmed
gain once procedure overlap is controlled.

\subsection{Relationship to prior work}
\label{sec:prior}
Two earlier papers by the authors are the direct starting point. \cite{b31}
introduced the retrieval and generation pipeline for first-reply generation;
\cite{b32} added LLM-judge feedback as a retrieval modifier and reported an
exploratory, conditioned-only analysis on 250 leave-one-out queries with a
different judge, in which no always-on method was significant and retrieval-score
gates were indistinguishable from always-on. The present article re-uses the
pipeline and the legacy lift formula as fixed components. Its contribution is
different: it treats the data regime as the experimental variable, audits the
earlier split, restricts retrieval, feedback and priors to an allowed history,
isolates the components of the calibration repair, adds a reply-similarity
baseline that uses no LLM judge and a matched-retriever comparison, and reports a blinded expert
assessment of procedural correctness. Table~\ref{tab:prior} compares the studies
side by side. No text, tables or figures are reused; the overlap is limited to the
cited background, and the connection is disclosed here as required for expanded
conference versions.

\Table[!btp](width=\textwidth){Relationship to the authors' prior conference
papers and to the earlier evaluation of the same pipeline. This article shares the
corpus and the base pipeline; the data-regime control, split audit, controlled
repair, reply baseline, matched retrievers, statistical sensitivity analysis and
expert evaluation are new.\label{tab:prior}}{
\scriptsize
\begin{tabularx}{\textwidth}{@{}p{0.15\textwidth}>{\raggedright\arraybackslash}X>{\raggedright\arraybackslash}X>{\raggedright\arraybackslash}X>{\raggedright\arraybackslash}X@{}}
\toprule
 & SIKDD 2025 & SIKDD 2026 & Earlier evaluation & This article \\
\midrule
Role & Retrieval + generation pipeline & Exploratory feedback-aware retrieval & Protocol as experimental variable & Data regime and measurement validity \\
Corpus & Enterprise IT tickets & 1,596 tickets, 37 teams & 1,595 tickets, 37 teams, 17 classes & Same 1,595; canonical IDs unchanged \\
Evaluation queries & --- & 250 leave-one-out & 878/319/398 + 5 seeds; 390 ``disjoint'' eval & Same legacy splits; history-only 398/398/427 \\
Judge & --- & gpt-4o-mini, conditioned only & gpt-5.6-luna, conditioned and blind & Same two protocols, restricted to history \\
Feedback pairs & --- & 102,722 conditioned & 175,600 per protocol & 94,913 / 52,321 / 58,826 admissible \\
Lift & --- & Laplace, $m=0.80$, cap $0.20$ & Legacy Laplace as object of study & Tuned Laplace, EB, scaling, scope, backoff \\
Procedure-disjoint & No & No & Claimed, 390 eval; not actually disjoint & Versioned groups; duplicate and prefix panels; 5 seeds \\
History-only evidence & No & No & No & Yes: index, BM25, feedback rows, priors \\
Reply-similarity baseline & No & No & No & Exact and BGE reply similarity \\
Matched retrievers & No & No & Confounded lift settings & One configuration; dense, BM25, hybrid \\
Human evaluation & No & No & No & 105 blinded items; E1 complete, E2 pending \\
Statistics & --- & Bootstrap CIs; no always-on significance & Dev-frozen thresholds; 39 runs; no multiplicity correction & Pre-specified families, Holm, MDE, exact McNemar \\
Main finding & Pipeline works & Conditioned feedback helps weak baselines; gates fail & Same reranking helps conditioned, harms blind & Harm is overlap-driven; calibrated gain small and metric-dependent; expert and cosine diverge \\
\bottomrule
\end{tabularx}}

\subsection{Gap}
Across these literatures, the information available to the feedback source is
rarely separated from the mechanism that converts evidence into scores, and the
overlap between evaluation tickets and the historical base is rarely measured.
This makes ``feedback helps'' and ``feedback hurts'' results hard to reconcile.
It is also rare to report where a control mechanism removes harm and where it
merely removes no-ops. This paper separates the two axes experimentally and asks
whether the gain survives a strict data boundary. A practical consequence,
developed in Section~\ref{sec:repair}, is that calibration and control should be
compared as tested configurations rather than assumed to be independent
improvements.

\subsection{RAG evaluation and procedural correctness}
RAG evaluation suites score faithfulness, answer relevance and context relevance
against retrieved evidence \cite{b40,b41,b42}, and retrieval benchmarks provide
zero-shot comparisons across domains \cite{b43}. Those measures answer whether an
answer is supported by its context; they do not answer whether the answer describes
the procedure that resolves the ticket. A reply can be faithful to a retrieved
passage and still recommend the wrong form, redirect or action. Our expert
assessment is a reminder that the two questions are different. The saturation
diagnostic we develop---the share of candidates pinned at the lift cap and the
share of exact-zero judge scores---is cheap to compute on any feedback source and
can be used as a pre-deployment check even when no expert is available.
